% Casey/W33 RH next-five insertion: orbit positivity, prime obstruction,
% infinite-operator falsifier, de Branges boundary, and norm-11 elliptic bridge.

\section{The Casey--W33 Reflection Cocycle and the Norm-$11$ Elliptic Packet}
\label{sec:casey_w33_weil_elliptic}

\subsection{Orbit-level positive quadratic form}

Let
\[
 \rho=\frac12+\delta+i\gamma,
 \qquad
 \rho^{\star}=1-\overline\rho=\frac12-\delta+i\gamma,
\]
and put $\sigma=\frac12+a$ with $a>|\delta|$. The reflection-cocycle
current is the logarithmic phase current of the ratio
\[
 R_\rho(s)=\frac{s-\rho}{s-\rho^{\star}}.
\]
After translating $y=t-\gamma$, its squared boundary energy is
\begin{equation}
 \mathcal E_{a,\delta}
 =\int_{\mathbb R}|A_{a,\delta}(y)|^2\,dy
 =\frac{\pi\delta^2}{a(a^2-\delta^2)}.
 \label{eq:casey_energy_next5}
\end{equation}
It has the exact Laplace-space representation
\begin{equation}
 \frac{\mathcal E_{a,\delta}}{\pi}
 =\int_0^\infty
 \left(e^{-(a-\delta)x}-e^{-(a+\delta)x}\right)^2\,dx.
 \label{eq:casey_laplace_norm}
\end{equation}
Thus the repaired torque is a positive Cauchy-kernel norm. Any convergent
nonnegative sum of such orbit energies vanishes if and only if every included
orbit lies on $\Re(s)=1/2$.

\paragraph{Claim boundary.}
Equations~\eqref{eq:casey_energy_next5}--\eqref{eq:casey_laplace_norm} give an
orbit-level RH detector. They do not prove that the corresponding classical
sum over all zeta-zero orbits converges or vanishes.

\subsection{Prime-power second difference and the boundary obstruction}

Sampling~\eqref{eq:casey_laplace_norm} at $x=\log n$ gives the positive
von Mangoldt sum
\begin{equation}
 \mathcal P(a,\delta)
 =\sum_{n\ge2}\Lambda(n)n^{-2a}
       \left(n^\delta-n^{-\delta}\right)^2.
\end{equation}
Writing $L(s)=-\zeta'(s)/\zeta(s)$, absolute convergence for
$a-|\delta|>1/2$ yields
\begin{equation}
 \mathcal P(a,\delta)
 =L(2(a-\delta))-2L(2a)+L(2(a+\delta)).
 \label{eq:prime_second_difference}
\end{equation}
Since $a=\sigma-1/2$, the positive sum converges only for
$\sigma>1+|\delta|$. In particular, the proposed boundary $\sigma=1$ is
outside the absolute-convergence region for every nonzero defect. A global
classical identity must therefore use the regularized Weil explicit formula,
including its archimedean and endpoint terms; the raw positive prime sum is
not sufficient.

\subsection{Positive infinite tower and its eighth-moment falsifier}

Let $\theta_j$ be the two self-adjoint W33 phase angles, with multiplicities
$24$ and $15$. A prime-indexed positive trace model with weight
$\Lambda(n)(\log n)^2n^{-s}$ and inverse local ordinate
$\log(n)/\theta_j$ has moments
\begin{equation}
 M_{2k}(s)=T_{2k}\frac{d^{2k+2}}{ds^{2k+2}}
          \left[-\frac{\zeta'(s)}{\zeta(s)}\right],
 \qquad s>1.
\end{equation}
A positive mixture of the $s=2$ and $s=3$ channels, followed by one overall
scale and trace normalization, can interpolate the first three even Taylor
moments of $\log\xi(1/2+z)$. This is a three-parameter interpolation. The
unfitted eighth moment differs from the classical value by
\[
 47.5186498\ldots\%.
\]
Hence this two-channel tower is falsified as a classical Hilbert--Polya
operator.

\subsection{Why generic de Branges positivity is insufficient}

For every $\delta\in\mathbb R$ and $\gamma>0$, the polynomial
\[
 E_{\delta,\gamma}(z)=(z+i\gamma)^2-\delta^2
\]
is Hermite--Biehler, because for $y>0$,
\begin{equation}
 |E_{\delta,\gamma}(x+iy)|^2
 -|E_{\delta,\gamma}^{\#}(x+iy)|^2
 =8\gamma y(\delta^2+\gamma^2+x^2+y^2)>0.
\end{equation}
This remains true for $\delta\ne0$. Therefore the existence of an
orbit-adapted de Branges space does not force the critical line. A classical
proof would need one fixed entire Hermite--Biehler function canonically tied
to $\xi$, with global kernel positivity and no orbit-dependent retuning.

\subsection{Exact local Hasse--Weil realization at $p=11$}

Consider the elliptic curves over $\mathbb Q$
\[
 E_2: y^2=x^3+x-1,
 \qquad
 E_{-4}: y^2=x^3+x+2.
\]
Both have good reduction at $11$, and direct counting gives
\[
 \#E_2(\mathbb F_{11})=10,
 \qquad
 \#E_{-4}(\mathbb F_{11})=16.
\]
Thus their Frobenius traces are $a_{11}=2$ and $a_{11}=-4$, respectively,
and their local polynomials are
\[
 P_{11}(E_2,u)=1-2u+11u^2,
 \qquad
 P_{11}(E_{-4},u)=1+4u+11u^2.
\]
Consequently the nontrivial W33 Ihara determinant factors exactly as
\begin{equation}
 Z_{W,\mathrm{nt}}(u)^{-1}
 =P_{11}(E_2,u)^{24}P_{11}(E_{-4},u)^{15}.
 \label{eq:w33_elliptic_local_factor}
\end{equation}
The Frobenius discriminants are $-40$ and $-28$, producing
$\mathbb Q(\sqrt{-10})$ and $\mathbb Q(\sqrt{-7})$, exactly the two
quadratic fields of the W33 nontrivial pole coordinates.

If $T_n=\alpha^n+\beta^n$ for either Frobenius pair, then
\begin{equation}
 T_0=2,
 \qquad
 T_1=a_{11},
 \qquad
 T_n=a_{11}T_{n-1}-11T_{n-2}.
\end{equation}
This is simultaneously the Hashimoto sector recurrence, while
\[
 \#E(\mathbb F_{11^n})=11^n+1-T_n.
\]
Therefore the two W33 nonbacktracking sectors are exactly two elliptic
Frobenius sectors at the local prime $11$, at every extension degree.

\paragraph{Global obstruction.}
Equation~\eqref{eq:w33_elliptic_local_factor} is an exact local theorem. It
does not construct an adelic object whose full $L$-function is the completed
Riemann zeta function. A global transfer must explain the multiplicities
$24$ and $15$, produce compatible factors at all primes, and match the
archimedean completed factor.
